\chapter{Conclusión y trabajo futuro}\label{ch:conclusion}


Un dato revelador es que el WER sobre el conjunto de entrenamiento se sitúa en torno al $7\%$, frente al $27\%$ obtenido en dev y test.

Esta brecha de aproximadamente $20$ puntos porcentuales es una señal clara de sobreajuste: el modelo memoriza con alta precisión las secuencias de glosas del corpus de entrenamiento pero generaliza de forma limitada a muestras no vistas.

Este fenómeno es esperable dado el tamaño reducido del corpus ($7.096$ muestras de entrenamiento) y la complejidad de la arquitectura.

El hecho de que ninguna de las técnicas de aumentación logre cerrar significativamente esta brecha refuerza la hipótesis de que el problema no es trivialmente soluble mediante perturbaciones sobre los keypoints, y apunta a la necesidad de datos adicionales o estrategias de regularización más agresivas como líneas de trabajo futuro.
